% !TeX root = ../main.tex
\section{Task Definition and Evaluation}
\label{sec:task}

\subsection{Outputs}
A complete tablature for one guitar part contains (i) note events with onset,
offset and pitch, (ii) a string and fret for each note, given a tuning and
possibly a capo, (iii) a rhythmic notation: beats, bars, time signatures and a
quantised duration for each note, and (iv) technique annotations per note
(bend height and release, vibrato, slide, hammer-on and pull-off, palm mute,
harmonics, dead notes). When the recording contains several guitars, each
part must also be separated into its own track. Published systems address
subsets of these outputs; very few address (iii), and none address all of
them from a band recording.

\subsection{Metrics}
\emph{Note-level metrics} follow the \texttt{mir\_eval} conventions used for
piano~\cite{hawthorne2018onsets}: a predicted note matches a reference note if
the pitches agree (within 50 cents) and the onsets lie within 50\,ms; the
stricter variant also requires the offset within
$\max(50\,\text{ms}, 0.2\,d)$ of the reference, $d$ being its duration.
Precision, recall and F1 follow from a one-to-one matching. We write F50 for
the onset-only F1. MT3~\cite{gardner2022mt3} noted that multi-instrument
transcription lacked consistent metrics; MuScriptor~\cite{rouard2026muscriptor}
reports instrument-agnostic onset, offset and frame F1 on its own test set.

\emph{Tablature metrics} add the string. The frame-level metrics of
TabCNN~\cite{wiggins2019tabcnn} (tablature precision, recall and F-measure,
and multi-pitch metrics) count a frame as correct when the right fret is
active on the right string. Note-level Tab F1, used by
TART~\cite{gupta2026tart}, requires a note match plus the same string and
fret. Symbolic fingering methods, which start from the reference notes, report
string or tablature accuracy against a human tablature, and some add
playability measures or user
studies~\cite{edwards2024miditotab,hamberger2025fretting}. Because several
fingerings of the same passage are idiomatic, agreement with one human
tablature is a lower bound on quality.

\emph{Rhythm metrics} are those of beat tracking: beat and downbeat F-measure
with a 70\,ms window, and continuity metrics such as CMLt, which require
consecutive correct beats~\cite{foscarin2024beatthis}. Systems that output
notated rhythm rarely report a metric for it.

\emph{Technique metrics} are per-class precision, recall and F1, either on
isolated notes or on notes of a transcription~\cite{su2019tent,gupta2026tart}.

\subsection{Comparability}
Reported numbers are hard to compare for three reasons. Test sets differ in
instrument (acoustic, classical, electric), in recording condition (direct
input, amplified, miked) and in whether the input is solo or mixed. Many
results are obtained on the same dataset the model was trained on, while the
use case is zero-shot. And labels carry systematic timing offsets: hexaphonic
pickups and manual annotations can be shifted by tens of milliseconds, which
is comparable to the 50\,ms tolerance, so the same model can score
differently depending on label alignment.
